% Part: second-order-logic
% Chapter: sol-and-set-theory
% Section: comparing-sets

\documentclass[../../../include/open-logic-section]{subfiles}

\begin{document}

\olfileid{sol}{set}{cmp}

\olsection{مقایسهٔ مجموعه‌ها}

\begin{prop}
!!{formula}~$\lforall[x][(X(x) \lif Y(x))]$ رابطهٔ زیرمجموعه‌بودن را
تعریف می‌کند؛ یعنی~$\Sat{M}{\lforall[x][(X(x) \lif Y(x))]}[s]$ اگر و
تنها اگر~$s(X) \subseteq s(Y)$.
\end{prop}

\begin{prop}
!!{formula}~$\lforall[x][(X(x) \liff Y(x))]$ رابطهٔ همانی بر مجموعه‌ها
را تعریف می‌کند؛ یعنی~$\Sat{M}{\lforall[x][(X(x) \liff Y(x))]}[s]$ اگر
و تنها اگر~$s(X) = s(Y)$.
\end{prop}

\begin{prop}
!!{formula}~$\lexists[x][X(x)]$ ویژگی ناتهی‌بودن را تعریف می‌کند؛ یعنی
$\Sat{M}{\lexists[x][X(x)]}[s]$ اگر و تنها اگر
$s(X) \neq \emptyset$.
\end{prop}

مجموعهٔ~$X$ از مجموعهٔ~$Y$ بزرگ‌تر نیست، یعنی~$\cardle{X}{Y}$، اگر و
تنها اگر تابعی~!!a{injective} چون~$f\colon X \to Y$ وجود داشته باشد. چون
می‌توانیم یک‌به‌یک‌بودن تابع و نیز قرارگرفتن ارزش‌های آن برای آرگومان‌های
$X$ در~$Y$ را بیان کنیم، می‌توانیم رابطهٔ بزرگ‌ترنبودن را بر
زیرمجموعه‌های دامنه نیز تعریف کنیم.

\begin{prop}
فرمول
\[
\lexists[u][(\lforall[x][(X(x) \lif Y(u(x)))] \land \lforall[x][\lforall[y][(\eq[u(x)][u(y)] \lif
      \eq[x][y])]])]
\]
رابطهٔ بزرگ‌ترنبودن را تعریف می‌کند.
\end{prop}

دو مجموعه هم‌اندازه، یا «هم‌شمار»، هستند، یعنی~$\cardeq{X}{Y}$، اگر و
تنها اگر تابعی~!!a{bijective} چون~$f\colon X \to Y$ وجود داشته باشد.

\begin{prop}
فرمول
\begin{multline*}
  \lexists[u][(\lforall[x][(X(x) \lif Y(u(x)))] \land {}\\
    \lforall[x][\lforall[y][((X(x) \land X(y) \land \eq[u(x)][u(y)]) \lif \eq[x][y])]]
  \land {} \\ \lforall[y][(Y(y) \lif \lexists[x][(X(x)
      \land \eq[y][u(x)])])])]
\end{multline*}
رابطهٔ هم‌شماربودن را تعریف می‌کند.
\end{prop}

این !!{formula}s را به‌ترتیب با~$X \subseteq Y$، $X = Y$،
$X \neq \emptyset$، $\cardle{X}{Y}$ و~$\cardeq{X}{Y}$ کوتاه می‌کنیم.
(این کار شاید اندکی گیج‌کننده باشد، زیرا هنگامی که به‌طور غیرصوری
دربارهٔ مجموعه‌های~$X$ و~$Y$ سخن می‌گوییم نیز همین نمادگذاری را به کار
می‌بریم؛ اما اینجا نمادگذاری، کوتاه‌نوشتِ !!{formula}sی در منطق مرتبهٔ
دوم است که متغیرهای رابطه‌ای یک‌موضعی~$X$ و~$Y$ را در بر دارند.)

\begin{prop}
!!{sentence}~$\lforall[X][\lforall[Y][((\cardle{X}{Y} \land
    \cardle{Y}{X}) \lif \cardeq{X}{Y})]]$ معتبر است.
\end{prop}

\begin{proof}
این !!{sentence} در !!a{structure}ی چون~$\Struct{M}$ ارضا می‌شود اگر
برای هر دو زیرمجموعه~$X \subseteq \Domain{M}$ و
$Y \subseteq \Domain{M}$، از~$\cardle{X}{Y}$ و~$\cardle{Y}{X}$ نتیجه
شود~$\cardeq{X}{Y}$. اما این نتیجه برای \emph{هر} دو مجموعهٔ~$X$ و~$Y$
برقرار است: این همان قضیهٔ شرودر--برنشتاین است.
\end{proof}


\end{document}
